\chapter{CLI and API Quick Reference}
\index{Azure CLI}\index{REST API}\index{quick reference}\index{T-SQL!reference}\index{KQL!reference}\index{PySpark!reference}%
This appendix collects the commands, endpoints, and idioms the book's labs use repeatedly, in
one lookup-friendly place. It is a reference, not a tutorial: each section points back to the
chapter that teaches the concept.

\section{Azure CLI Essentials}
\index{az CLI!essentials}%
\begin{lstlisting}[language=bash,caption={Azure CLI commands used across the labs.},label={lst:appendixb-az}]
az login                                  # authenticate (Ch. 0, Appendix A)
az account show                           # confirm subscription context
az group create -n rg-fabric-de -l eastus # resource group for lab infra

# Azure SQL free-tier source for the Week 9/11 labs
az sql server create -g rg-fabric-de -n sql-fabric-de-<suffix> -l eastus `
    -u sqladmin -p '<from-vault>'
az sql db create -g rg-fabric-de -s sql-fabric-de-<suffix> -n sourcedb `
    --edition GeneralPurpose --compute-model Serverless -f Gen5 -c 1

# Key Vault for Chapters 11, 19, 22 (secrets never in code)
az keyvault create -g rg-fabric-de -n kv-fabric-de-<suffix> -l eastus
az keyvault secret set --vault-name kv-fabric-de-<suffix> `
    -n eventstream-conn-str --value '<connection string>'
az keyvault secret show --vault-name kv-fabric-de-<suffix> `
    -n eventstream-conn-str --query value -o tsv

# Capacity hygiene (Ch. 5, 23): pause dev/test capacities off-hours
az fabric capacity resume -g rg-fabric-de -n fcapfabricde
az fabric capacity suspend -g rg-fabric-de -n fcapfabricde
\end{lstlisting}

\section{OneLake Paths and Endpoints}
\index{OneLake!paths}%
Every Fabric item with storage is addressable through the DFS and Blob endpoints. The two forms
appear throughout the book:

\begin{lstlisting}[language=bash,caption={OneLake addressing patterns.},label={lst:appendixb-onelake}]
# DFS endpoint (used by COPY INTO, ADLS-compatible tools)
https://onelake.dfs.fabric.microsoft.com/<workspace>/<item>/Files/<path>
https://onelake.dfs.fabric.microsoft.com/<workspace>/<item>/Tables/<table>

# ABFS form (used by Spark and storage SDKs)
abfss://<workspace>@onelake.dfs.fabric.microsoft.com/<item>.Lakehouse/Files/<path>
\end{lstlisting}

Right-click any folder in the lakehouse explorer to copy these paths exactly; a wrong segment is
the most common cause of failed \texttt{COPY INTO} loads (Chapter~5).

\section{Fabric REST API}
\index{Fabric REST API}%
The Fabric REST API automates workspace and item management---the operations plane beneath
Week~23's CI/CD. Pattern: acquire a token, call the API, poll long-running operations.

\begin{lstlisting}[language=bash,caption={Fabric REST API essentials.},label={lst:appendixb-rest}]
# Token for the Fabric API (service principal or az CLI context)
$token = az account get-access-token `
    --resource https://api.fabric.microsoft.com --query accessToken -o tsv

# List workspaces
curl -H "Authorization: Bearer $token" `
    https://api.fabric.microsoft.com/v1/workspaces

# Create a lakehouse in a workspace
curl -X POST -H "Authorization: Bearer $token" -H "Content-Type: application/json" `
    -d '{"displayName":"lh_lab","description":"Appendix B example"}' `
    https://api.fabric.microsoft.com/v1/workspaces/<workspace-id>/lakehouses

# Run a notebook job (long-running; poll the returned Location header)
curl -X POST -H "Authorization: Bearer $token" -H "Content-Type: application/json" `
    -d '{}' `
    "https://api.fabric.microsoft.com/v1/workspaces/<ws>/items/<notebook-id>/jobs/instances?jobType=RunNotebook"
\end{lstlisting}

\begin{pitfall}
The REST API authenticates as \emph{someone}: a service principal with least-privilege workspace
assignment, vaulted secret, and a rotation date. API tokens pasted into scripts reproduce
Chapter~19's key-in-notebook failure at the platform layer.
\end{pitfall}

\section{T-SQL Quick Reference}
\index{T-SQL!quick reference}%
The statements the warehousing chapters rely on, in one block:

\begin{lstlisting}[language=SQL,caption={T-SQL idioms used across Chapters 5--8 and 21.},label={lst:appendixb-tsql}]
-- Bulk load (Ch. 5)
COPY INTO dbo.stg_sales
FROM 'https://onelake.dfs.fabric.microsoft.com/<ws>/<lh>/Files/stage/sales/'
WITH ( FILE_TYPE = 'PARQUET' );

-- Idempotent manifest recording (Ch. 5, 9)
INSERT cfg.load_manifest (batch_path, rows_loaded, loaded_ts)
VALUES ('2026-07', @@ROWCOUNT, CURRENT_TIMESTAMP);

-- Cross-item query (Ch. 5-6): three-part name into a lakehouse endpoint
SELECT TOP 10 * FROM lh_gold.dbo.fact_sales;

-- Object security (Ch. 6, 21)
GRANT SELECT ON gold.v_sales_enriched TO [analysts@contoso.com];
DENY  SELECT ON gold.dim_customer (national_id) TO [analysts@contoso.com];

-- Row-level security (Ch. 6, 21)
CREATE SECURITY POLICY sec.sales_region_policy
ADD FILTER PREDICATE sec.fn_region_predicate(region)
ON gold.v_sales_enriched_base
WITH (STATE = ON);

-- Test as a restricted principal (Ch. 6, 21)
EXECUTE AS USER = 'analyst-us@contoso.com';
SELECT DISTINCT region FROM gold.v_sales_enriched;
REVERT;
\end{lstlisting}

\section{KQL Quick Reference}
\index{KQL!quick reference}%
\begin{lstlisting}[language=SQL,caption={KQL idioms used across Chapters 13--16.},label={lst:appendixb-kql}]
// Always bound by time (Ch. 14)
IoTTelemetry | where Timestamp > ago(1h);

// Dedup: latest arrival per producer-stamped id (Ch. 13-14)
IoTTelemetry | summarize arg_max(ingestion_time(), *) by EventId;

// Windowed aggregation (Ch. 14)
IoTTelemetry
| summarize avg(Temperature) by DeviceId, bin(Timestamp, 5m);

// Rolling average + anomalies (Ch. 14)
IoTTelemetry
| where Timestamp > ago(6h)
| make-series AvgTemp = avg(Temperature) default=0 on Timestamp step 5m by DeviceId
| extend RollingAvg = series_fir(AvgTemp, repeat(1, 5))
| extend Anomalies = series_decompose_anomalies(RollingAvg, 1.5);

// Impossible travel (Ch. 16)
Txn
| summarize arg_max(ingestion_time(), *) by txn_id
| order by card_id asc, ts asc
| serialize
| extend p_city = prev(city), p_ts = prev(ts)
| where city != p_city and ts - p_ts <= 10m;
\end{lstlisting}

\section{PySpark and Delta Quick Reference}
\index{PySpark!quick reference}\index{Delta Lake!quick reference}%
\begin{lstlisting}[language=Python,caption={PySpark/Delta idioms used throughout.},label={lst:appendixb-pyspark}]
# Read a managed table; stamp lineage columns (Ch. 3)
df = (spark.read.table("bronze.orders")
      .withColumn("ingested_at", F.current_timestamp()))

# Idempotent partition-scoped write (Ch. 9, 17, 20)
(df.write.format("delta").mode("overwrite")
   .option("replaceWhere", "as_of_date = '2026-07-31'")
   .saveAsTable("gold.features_store_daily"))

# Delta maintenance (Ch. 4)
spark.sql("OPTIMIZE silver_orders VORDER")
spark.sql("DESCRIBE HISTORY silver_orders")

# Time travel (Ch. 4, 22)
prev = (spark.read.format("delta").option("versionAsOf", 12)
        .table("silver_sales_orders_training"))

# Schema evolution (Ch. 9): default enforcement rejects; opt in deliberately
(df.write.format("delta").mode("append")
   .option("mergeSchema", "true")
   .saveAsTable("silver.orders"))

# Registered model as UDF (Ch. 18, 20)
predict = mlflow.pyfunc.spark_udf(spark, "models:/churn_model/Production",
                                  result_type="double")
\end{lstlisting}

\section{DAX Quick Reference}
\index{DAX!quick reference}%
\begin{lstlisting}[language=SQL,caption={DAX idioms used across Chapters 7--8.},label={lst:appendixb-dax}]
Total Revenue = SUM ( fact_sales[net_amount] )

Order Count = DISTINCTCOUNT ( fact_sales[order_id] )

Average Order Value = DIVIDE ( [Total Revenue], [Order Count] )

Revenue YTD =
CALCULATE ( [Total Revenue], DATESYTD ( dim_date[full_date] ) )
\end{lstlisting}

\begin{tipbox}
Keep this appendix open during capstones. Every idiom here appears in at least one chapter's
validation checklist; when an acceptance criterion says ``idempotent'' or ``dedup,'' these are the
exact patterns it refers to.
\end{tipbox}
